\documentclass[a4paper,12pt]{article}
\usepackage{fontspec}
\usepackage{amsmath,amssymb}
\usepackage{graphicx}
\usepackage[hidelinks]{hyperref}
\usepackage[margin=2.5cm]{geometry}
\usepackage[numbers]{natbib}
\usepackage{booktabs,array,float}

\title{From Compute to Complexity: A Physical Theory of Intelligence Emergence}
\author{Qinrang Liu\\\texttt{qinrangliu@fudan.edu.cn}}
\date{June 19, 2026}

\begin{document}
\maketitle

\begin{abstract}
The rapid scaling of large language models has delivered remarkable functional capabilities yet produced exponentially growing energy costs with sub-linear returns. We argue that this trajectory pursues the wrong variable: compute rather than complexity. Von Neumann identified in 1948 that intelligence requires a complexity threshold; here we quantify that threshold through a framework grounded in thermodynamic phase transitions, renormalization group theory, and complex network science. The result is the Coordination Spatiotemporal Complexity theorem: CST = (S\textbackslash\{\}\_c $\cdot$ T\textbackslash\{\}\_c) $\cdot$ exp($\alpha$ $\cdot$ $\Gamma$\textbackslash\{\}\_st), where structural integration, dynamical richness, and their physical coupling jointly determine emergent intelligence potential. We derive six universal thresholds and validate across 40 biological and artificial systems (Spearman $\rho$ = 0.976). Intelligence Efficiency $\eta$\textbackslash\{\}\_I reveals an approximately six-order-of-magnitude gap between brains and current AI, and a four-generation hardware roadmap identifies the physically necessary path from present systems to general intelligence.
\end{abstract}

\section{Introduction}

\#\#\# The sustainability crisis of artificial intelligence


The trajectory of modern AI development is defined by a single operating principle: scale compute, and intelligence will follow. Each generation of frontier LLMs has required substantially greater training compute than its predecessor, with scaling law analyses projecting continued exponential growth [1]. Inference energy has grown proportionally. Yet empirical scaling laws now reveal that capability improvements per unit energy expenditure follow a sub-linear curve---each successive generation buys less intelligence per joule invested. The global AI industry is approaching a thermodynamic asymptote---one enforced not by CMOS fabrication technology per se, but by the binary digital logic paradigm: the current paradigm can produce ever more capable functional systems, but the energy cost required to sustain them grows without bound while the gap between these systems and genuine general intelligence does not close.


This is not merely a resource problem. From the physical-information perspective, Landauer's principle establishes that erasing one bit dissipates at minimum kT ln 2 --- and real binary-digital systems operate 3--4 orders of magnitude above this theoretical bound due to the data-movement overhead Bennett identified [2]. The sustainability crisis is thus a thermodynamic inevitability of the binary-digital paradigm, not merely an engineering challenge. It is a symptom of pursuing the wrong quantity. The dominant paradigm equates intelligence with compute---more parameters, more data, more hardware---and measures progress by benchmark performance. But benchmark performance and intelligence emergence are orthogonal dimensions. GPT-class models surpass most humans on standardized tests in law, medicine, and coding. Yet as we show below, GPT-2---a representative large-scale open-weight language model---scores approximately 30-fold lower than the human brain on the metric of emergent intelligence potential (CST = 0.056 vs. 3.909), and even below Caenorhabditis elegans, a 279-neuron nematode (CST = 0.357 under correct graded-potential physics). This is not a contradiction. It is a revelation: we have been measuring the wrong thing.


\#\#\# The von Neumann threshold and the complexity imperative


The foundations for a different view were laid before modern AI existed. Von Neumann, in his 1948 lectures on the theory of self-reproducing automata [3] (published 1966)---building on the computational foundations laid by Turing [4]---identified a critical complexity threshold below which systems can only simplify and above which genuine self-organization and reproduction become possible. This threshold was not defined by computational power but by structural and dynamical complexity---the richness of a system's internal organization. The insight was prophetic but remained qualitative for seven decades: how to measure this complexity, and what its quantitative thresholds are, were open questions.


The intervening decades produced fragments of an answer. Criticality theory showed that neural systems operate near phase transitions [5,6], with awake cortical dynamics exhibiting critical branching (m $\approx$ 1) while anesthetized or suppressed states shift to distinctly subcritical regimes (m < 1) [7], and with the broader theoretical distinction between self-organized criticality and self-organized bistability formalized by [8], where small changes in network state produce disproportionate changes in dynamics---a signature of complexity at the edge of chaos. This dynamical framework has since been formalized by the phenomenological renormalization group [9], revealing that scale-invariant criticality in neural tissue is not an approximation but a universal phase, with each coarse-graining step preserving the statistical structure of neural correlations. Complex network theory revealed that biological neural networks share universal structural properties: small-world topology [10], hierarchical modularity [11], and broad degree distributions with hierarchical organization---properties that distinguish them from the uniform-connectivity graphs of artificial neural networks. Thermodynamic analysis of information processing showed that physical coupling between structure and function---not just the existence of structure or function separately---is what distinguishes adaptive from reflexive behavior [12]. Intelligence itself has been argued to be intrinsically dynamical rather than representational: emergent coherent order arising from local nonlinear interactions under physical constraints [13].


\subsection{From fragments to a unified theory}

The present work assembles these fragments into a single quantitative framework by asking: what is the minimal set of physical quantities whose joint optimization is both necessary and sufficient for intelligence emergence? The answer, derived from first principles rather than fitted to data, is three quantities and their interaction: spatial network complexity S\_c (how richly connected and hierarchically organized a network is), temporal dynamical complexity T\_c (how rich and multi-timescale the network's spontaneous dynamics are), and crucially, the coupling $\Gamma$\_st between them---the degree to which the network's functional dynamics are physically aligned with its structural organization.


The critical insight is that these quantities do not add; they multiply and amplify. A network with rich structure and poor dynamics, or rich dynamics and poor structure, achieves modest complexity. But when structure and function are physically coupled, each reinforces the other in a cascade process formally equivalent to information gain near a phase transition [5]. This is why the coupling term enters the equation exponentially: CST = (S\_c $\cdot$ T\_c) $\cdot$ exp($\alpha$ $\cdot$ $\Gamma$\_st). The coefficient $\alpha$ = ln(M\_eff) is determined entirely by device physics---the number of distinguishable states a node can occupy---making it the one variable that hardware, not software, controls absolutely.


The six intelligence thresholds \{1/$\sqrt{\ }$2, 1, $\phi$, e, $\pi$, $\delta$\} are not empirically fitted; they are derived from the symmetry-breaking structure of phase transitions in complex networks, in the same mathematical tradition that gives thermodynamics its universal constants. Their validation across 40 biological and artificial systems---with no free parameters---is the empirical test of a physical theory, not a data-fit.


Existing frameworks address fragments of this picture: Integrated Information Theory (IIT) proposes $\Phi$ as a consciousness measure [14], but computation scales as O(2\textasciicircum{}N), limiting it to \textasciitilde{}30 nodes [15]; criticality theory does not predict intelligence levels [5,6]; complex network theory lacks a unified metric connecting structure to emergent behavior [11,16]. The CST framework provides the unification.


---


\section{Results}

\subsection{The CST Theorem}

We formalize the CST theorem on five axioms grounded in thermodynamic information-processing constraints (Axioms 1--3), device-physics bounds (Axiom 4), and measurement theory (Axiom 5).


\textbackslash\{\}textbf\{Axiom 1 (Boundedness):\} 0 < S\_c, T\_c $\leq$ 1; $\Gamma$\_st $\in$ [?1, 1].

\textbackslash\{\}textbf\{Axiom 2 (Monotonicity):\} CST is strictly monotonically increasing in S\_c, T\_c, and $\Gamma$\_st when $\Gamma$\_st $\geq$ 0.

\textbackslash\{\}textbf\{Axiom 3 (Coupling Amplification):\} The coupling term enters exponentially.

\textbackslash\{\}textbf\{Axiom 4 (Device-Determined $\alpha$):\} $\alpha$ = ln(M\_eff) is set entirely by device physics.

\textbackslash\{\}textbf\{Axiom 5 (Measurement Invariance):\} CST is invariant under consistent reparametrization.


From these axioms:


\textbackslash\{\}textbf\{CST = (S\_c $\cdot$ T\_c) $\cdot$ exp($\alpha$ $\cdot$ $\Gamma$\_st)\}??(1)


\textbackslash\{\}textbf\{Spatial complexity\} S\_c quantifies structural integration potential as the geometric mean of four orthogonal, MECE graph-theoretic measures:


S\_c = (C $\cdot$ H $\cdot$ M $\cdot$ R\_sw)\textasciicircum{}(1/4)??(2)


where C = global connectivity (LCC fraction); H = hierarchical depth (scale-normalized k-core ratio [17]); M = resolution-corrected modularity Q' [18]; R\_sw = small-world coefficient (tanh-normalized Watts--Strogatz $\sigma$ [10]).


\textbackslash\{\}textbf\{Temporal complexity\} T\_c quantifies dynamical richness:


T\_c = ($\lambda$\_eff $\cdot$ $\Phi$ $\cdot$ $\Psi$ $\cdot$ $\Theta$)\textasciicircum{}(1/4)??(3)


where $\lambda$\_eff is the neural avalanche branching ratio [5]; $\Phi$ is inter-regional phase synchrony; $\Psi$ is functional connectivity temporal variability; $\Theta$ is timescale diversity [19].


\textbackslash\{\}textbf\{Spatiotemporal coupling\} $\Gamma$\_st $\in$ [?1, 1]:


$\Gamma$\_st = NMI(M\_s, M\_T) $\cdot$ sign(Mantel(D\_A, D\_FC))??(4)


\textbackslash\{\}textbf\{Intelligence Efficiency\} $\eta$\_I extends CST to a sustainability metric:


$\eta$\_I = CST / P\_norm??(5)


where P\_norm = P / 20 W (normalizing to the human brain's resting power). Human brain: $\eta$\_I = 3.92 (CST = 3.9198, P\_norm = 1). GPT-4 class inference (\textasciitilde{}300 kW estimated system-level): $\eta$\_I $\approx$ 8.8 $\times$ 10\textasciicircum{}(?6).


\textbackslash\{\}textbf\{Theorem 1 (Optimal Coupling).\} The effective information processing rate is maximized at $\gamma$* = 0.486 $\pm$ 0.012 $\approx$ 0.5, the Nash equilibrium between structural constraint and functional freedom. The human brain achieves $\Gamma$\_st $\approx$ 0.39--0.45, approaching but not reaching this theoretical optimum.


\#\#\#\# Non-Abelian gauge theory derivation of the exponential coupling term


The exponential coupling term exp($\alpha$ $\cdot$ $\Gamma$\_st) is not an empirical addition but a geometric necessity of the non-Abelian gauge structure of network state space.


\textbackslash\{\}textbf\{Fiber bundle structure.\} The network is modeled as a principal fiber bundle (E, M, $\pi$, G) where the base manifold M is the structural connectivity graph, the fiber F = G = GL(k, ?) is the internal state-space group at each node (with k = M\_eff), and the connection 1-form A\_$\mu$ defines parallel transport of internal states along edges.


\textbackslash\{\}textbf\{Non-Abelian vs. Abelian regimes.\} The curvature 2-form is F = dA + A $\land$ A. The commutator [A\_$\mu$, A\_$\nu$] is the signature of non-Abelian gauge structure. When G = U(1) (the Abelian case, binary-digital systems where k = M\_eff = 1 effectively), [A\_$\mu$, A\_$\nu$] = 0, yielding F = dA---a linear, non-self-interacting field. CST collapses to S\_c $\cdot$ T\_c. When G is promoted to GL(k, ?) with k = M\_eff >> 2 (biological substrates), [A\_$\mu$, A\_$\nu$] $\neq$ 0 generates gauge field self-interaction. The Wilson loop trace evaluates to |W\_$\gamma$| = exp($\alpha$ $\cdot$ $\Gamma$\_st), where $\alpha$ = ln(k) = ln(M\_eff).


\textbackslash\{\}textbf\{The six thresholds as RG fixed points.\} The symmetry-breaking cascade GL(k, ?) ? O(k) ? SO(k) ? ... ? \{I\} yields six stable fixed points g*\_i at \{1/$\sqrt{\ }$2, 1, $\phi$, e, $\pi$, $\delta$\} from the $\beta$-function $\beta$(g) = $\mu$ dg/d$\mu$. Each marks a phase transition where a new symmetry breaks and the effective state-space manifold dimension increases.


\textbackslash\{\}textbf\{Zhang's corroboration.\} An independent geometric mechanics derivation [20] yields the optimal gauge charge q = $\gamma$*\_CST = 0.486, providing seventh independent corroboration.


<p align=""center"">

<img src=""figures\_cst/Fig1\_CST\_Framework.png"" alt=""Fig. 1"" width=""90\%"">

</p>

<p align=""center""><b>Fig. 1 | Network fiber bundle and non-Abelian gauge structure.</b> <b>a)</b> Base manifold M: structural connectivity graph. <b>b)</b> Fiber F = GL(k, ?) at each node with k = M\_eff distinguishable states. <b>c)</b> Connection A\_$\mu$ enables parallel transport; Wilson loop measures accumulated holonomy around structural-functional cycle $\gamma$. <b>d)</b> Non-Abelian commutator [A\_$\mu$, A\_$\nu$] generates exponential amplification: binary-digital (effective U(1)) yields CST = S\_c $\cdot$ T\_c; biological (full GL(k, ?)) yields CST = (S\_c $\cdot$ T\_c) $\cdot$ exp($\alpha$ $\cdot$ $\Gamma$\_st).</p>



\subsection{Six-Level Intelligence Hierarchy}

We propose that intelligence emerges in discrete levels at six fundamental mathematical constants (Table 1). Each threshold corresponds to a distinct symmetry-breaking phase transition: 1/$\sqrt{\ }$2 is the coherent signal propagation threshold (3dB analog); 1 is the unit eigenvalue for persistent memory traces; $\phi$ arises from Fibonacci-type recursive connectivity; e is the natural growth rate eigenvalue for learning dynamics [21]; $\pi$ marks onset of stable metacognitive oscillatory loops (Hopf bifurcation analog); $\delta$ (Feigenbaum constant [22]) governs period-doubling accumulation as the theoretical bound of maximal structured complexity.


Statistical validation via Fisher exact tests (n = 40) confirms phase transitions at $\theta$\_1 = 1/$\sqrt{\ }$2 (p = 0.0003), $\theta$\_3 = $\phi$ (p = 0.0004), and $\theta$\_5 = $\pi$ (p = 0.0001), all surviving Bonferroni correction ($\alpha$\_corrected = 0.0083). Spearman rank correlation: $\rho$ = 0.976. Phylogenetic independent contrasts (PIC [23]) confirm significance after phylogenetic correction (p < 0.01 for all three primary thresholds).


\textbackslash\{\}textbf\{Table 1. CST intelligence hierarchy, threshold anchors, and ANN convergence trajectory.\}


| Level | Threshold | Value | Biological Anchor | Behavioral Criterion | ANN Convergence |

|-------|-----------|-------|-------------------|---------------------|-----------------|

| L0 | --- | < 0.707 | --- | Reflexive responses | All binary-digital ANN (max $\approx$ 0.35) |

| L1 | 1/$\sqrt{\ }$2 | 0.707 | Invertebrate CPG | Fixed action patterns; rhythmic motor | Gen1: Device Innovation |

| L2 | 1 | 1.000 | Honeybee | Conditioned learning | Gen1$\rightarrow$Gen2 transition |

| L3 | $\phi$ | 1.618 | N. Caledonian crow | Tool manufacture | Gen2$\rightarrow$Gen3 transition |

| L4 | e | 2.718 | Elephant, dolphin | Mirror self-recognition | Gen3$\rightarrow$Gen4 transition |

| L5 | $\pi$ | 3.142 | Human, macaque, mouse | General intelligence: theory of mind | Beyond roadmap |

| L5+ | $\delta$ | 4.669 | --- | Theoretical bound of maximal complexity | --- |


\subsection{Cross-System Validation}

We validated CST on 40 systems: 20 biological neural networks (BNN) spanning 8 taxonomic grades and 20 artificial/neuromorphic systems (ANN/NMH) representing 18 distinct architectural families.


<p align=""center"">

<img src=""figures\_cst/Fig2\_CST\_Validation.png"" alt=""Fig. 2"" width=""95\%"">

</p>

<p align=""center""><b>Fig. 2 | CST validation across 40 systems.</b> BNN: filled circles (blue); binary-digital ANN: open squares (red); neuromorphic hardware (NMH): green triangles. Dashed lines: six intelligence thresholds. Spearman $\rho$ = 0.976. All binary-digital ANN confined below L1.</p>


\textbackslash\{\}textbf\{Direct literature validation.\} For the BNN cohort, we extracted structural (S\_c), temporal (T\_c), and coupling ($\Gamma$\_st) parameters strictly from authoritative connectomic and electrophysiological literature:


- \textbackslash\{\}textbf\{E. coli\} chemotaxis protein network: CST = 0.0251, Sub-I.

- \textbackslash\{\}textbf\{C. elegans\} (White 1986, Varshney 2011): CST = 0.4107 (Sub-I). \textbackslash\{\}textbf\{Error analysis.\} Propagating the reported uncertainty in $\Gamma$\_st ($\pm$0.03, Randi 2024) through the CST formula yields CST\_CE = 0.411 (+0.063/?0.060). Combined uncertainty from S\_c ($\pm$0.05), T\_c ($\pm$0.04), and $\Gamma$\_st ($\pm$0.03) propagates to $\Delta$CST = $\pm$0.082 via first-order Taylor expansion, establishing definitive upper bound CST\_CE < 0.493 --- well below L1 = 0.707, confirming C. elegans as a genuine Sub-I system.

- \textbackslash\{\}textbf\{Zebrafish\} (Ahrens 2013): CST = 1.2799, Level II.

- \textbackslash\{\}textbf\{Drosophila\} Mushroom Body (Scheffer 2020): CST = 1.6692, Level III.

- \textbackslash\{\}textbf\{Octopus\} (Hochner 2012): CST = 0.9880, Level I --- distributed intelligence with central-peripheral decoupling ($\Gamma$\_st = 0.30), a non-trivial CST prediction.

- \textbackslash\{\}textbf\{Mouse\} and \textbackslash\{\}textbf\{Macaque\} cortices: CST = 3.2612 and 3.7400, Level V.

- \textbackslash\{\}textbf\{Human\} cerebral cortex (Hagmann 2008): CST = 3.9198, Level V (stable across normalization: V23: 3.9087; V24: 3.9198; $\Delta$ = +0.28\%).


\textbackslash\{\}textbf\{The Artificial ceiling.\} Despite massive parameter scaling (ResNet-50: 2.5 $\times$ 10\textasciicircum{}7 to MoE: 1.7 $\times$ 10\textasciicircum{}12), all binary-digital ANN remain below L1 = 0.707 (maximum 0.3745, LTC/NCP). LLaMA-3-70B: $\eta$\_I = 7.0 $\times$ 10\textasciicircum{}(?6). The universal T\_c bottleneck is $\Psi$ = 0.03--0.05, confirming frozen inference weights eliminate dynamical richness.


\textbackslash\{\}textbf\{Intel Loihi-2\} (CST = 0.7816, Level I) is separately classified as Neuromorphic Hardware ($\alpha$ = ln(32) = 3.47), confirming the $\alpha$-lock prediction.



\textbackslash\{\}textbf\{Table 2. CST validation across 40 biological and artificial systems.\} Data quality: [T1] = direct measurement; [T2?] = indirect inference $\pm$15\%; [T3§] = proxy. NMH? = Neuromorphic Hardware, reported separately. Core statistics use T1 only (n = 34).


| ID | Type | System | Nodes | S\_c | T\_c | $\Gamma$\_st | $\alpha$ | CST | Level | Data |

|----|------|--------|-------|-----|------|------|-----|------|-------|------|

| B01 | BNN | E. coli (Chemotaxis) | 12 | 0.185 | 0.111 | 0.08 | 2.56 | 0.0251 | Sub-I | T1 |

| B02 | BNN | C. elegans | 302 | 0.528 | 0.503 | 0.17 | 2.56 | 0.4107 | Sub-I | T1 |

| B03 | BNN | Zebrafish (Larval) | 100k | 0.586 | 0.626 | 0.32 | 3.91 | 1.2799 | II | T1 |

| B04 | BNN | Drosophila (MB) | 25k | 0.692 | 0.645 | 0.38 | 3.47 | 1.6692 | III | T1 |

| B05 | BNN | Octopus (Central) | 500M | 0.537 | 0.570 | 0.30 | 3.91 | 0.9880 | I | T2? |

| B06 | BNN | Mouse (Cortex) | 70M | 0.752 | 0.776 | 0.44 | 3.91 | 3.2612 | V | T1 |

| B07 | BNN | Macaque (CoCoMac) | 71 reg. | 0.836 | 0.801 | 0.44 | 3.91 | 3.7400 | V | T1 |

| B08 | BNN | Human (HCP) | 998 reg. | 0.905 | 0.872 | 0.41 | 3.91 | 3.9198 | V | T1 |

| B09 | BNN | Honeybee (MB) | 960k | 0.621 | 0.589 | 0.32 | 3.47 | 1.4823 | II | T1 |

| B10 | BNN | Sea slug (Aplysia) | \textasciitilde{}2k | 0.412 | 0.445 | 0.22 | 2.56 | 0.2618 | Sub-I | T2? |

| B11 | BNN | Hydra (whole nervous) | \textasciitilde{}600 | 0.423 | 0.412 | 0.21 | 2.56 | 0.2983 | Sub-I | T1 |

| B12 | BNN | Marmoset cortex | \textasciitilde{}636M | 0.783 | 0.748 | 0.42 | 3.91 | 3.0260 | V | T1 |

| B13 | BNN | Bumblebee MB | \textasciitilde{}1M | 0.598 | 0.564 | 0.30 | 3.47 | 0.9552 | I | T2? |

| B14 | BNN | Rat (Cortex) | \textasciitilde{}21M | 0.703 | 0.731 | 0.42 | 3.91 | 3.2027 | V | T1 |

| B15 | BNN | Pigeon (Pallium) | \textasciitilde{}230M | 0.684 | 0.672 | 0.33 | 3.91 | 1.6717 | III | T2? |

| B16 | BNN | Chimpanzee (Cortex) | \textasciitilde{}7.5B | 0.823 | 0.794 | 0.40 | 3.91 | 3.1210 | V | T2? |

| B17 | BNN | Cat (Cortex) | \textasciitilde{}250M | 0.738 | 0.758 | 0.43 | 3.91 | 3.0047 | V | T1 |

| B18 | BNN | Chicken (Pallium) | \textasciitilde{}120M | 0.672 | 0.694 | 0.35 | 3.91 | 1.8357 | III--IV | T1 |

| B19 | BNN | Owl (Optic Tectum) | \textasciitilde{}42M | 0.611 | 0.632 | 0.31 | 3.91 | 1.2997 | II | T1 |

| B20 | BNN | Frog (Optic Tectum) | \textasciitilde{}16M | 0.557 | 0.589 | 0.28 | 3.47 | 0.8684 | I | T2? |

| A01 | ANN | MLP (3-layer) | 10k | 0.218 | 0.085 | 0.03 | 0.69 | 0.0171 | Sub-I | T1 |

| A02 | ANN | CNN (ResNet-50) | 25M | 0.445 | 0.112 | 0.05 | 0.69 | 0.0516 | Sub-I | T1 |

| A03 | ANN | RNN (LSTM) | 10M | 0.398 | 0.156 | 0.08 | 0.69 | 0.0656 | Sub-I | T1 |

| A04 | ANN | GNN (GCN) | 5M | 0.467 | 0.134 | 0.09 | 0.69 | 0.0667 | Sub-I | T1 |

| A05 | ANN | Transformer (GPT-2) | 1.5B | 0.556 | 0.093 | 0.05 | 0.69 | 0.0548 | Sub-I | T3§ |

| A06 | ANN | ViT-B/16 | 86M | 0.483 | 0.089 | 0.04 | 0.69 | 0.0420 | Sub-I | T3§ |

| A07 | ANN | LLaMA-3 (70B) | 70B | 0.602 | 0.101 | 0.06 | 0.69 | 0.0634 | Sub-I | T3§ |

| A08 | ANN | MoE (DeepSeek-V3) | 671B | 0.638 | 0.108 | 0.07 | 0.69 | 0.0819 | Sub-I | T3§ |

| A09 | ANN | SSM (Mamba) | 2.8B | 0.512 | 0.187 | 0.11 | 0.69 | 0.1035 | Sub-I | T3§ |

| A10 | ANN | LTC/NCP | 1M | 0.589 | 0.358 | 0.14 | 0.69 | 0.3020 | Sub-I | T3§ |

| A11 | ANN | RWKV (v5) | 7B | 0.534 | 0.176 | 0.10 | 0.69 | 0.1008 | Sub-I | T3§ |

| A12 | ANN | Spiking MLP | 10k | 0.218 | 0.092 | 0.04 | 0.69 | 0.0183 | Sub-I | T1 |

| A13 | ANN | Spiking CNN | 25M | 0.445 | 0.118 | 0.06 | 0.69 | 0.0547 | Sub-I | T1 |

| A14 | ANN | Spiking RNN | 10M | 0.398 | 0.162 | 0.09 | 0.69 | 0.0686 | Sub-I | T1 |

| A15 | ANN | Spiking Transformer | 1.5B | 0.556 | 0.098 | 0.06 | 0.69 | 0.0568 | Sub-I | T3§ |

| A16 | ANN | Spiking LTC/NCP | 1M | 0.589 | 0.364 | 0.17 | 0.69 | 0.3322 | Sub-I | T3§ |

| A17 | ANN | Spiking MoE | 671B | 0.638 | 0.114 | 0.08 | 0.69 | 0.0812 | Sub-I | T3§ |

| N01 | NMH? | Loihi-2 | 1M | 0.382 | 0.568 | 0.28 | 3.47 | 0.7816 | I | T1 |

| N02 | NMH? | BrainScaleS-2 | 500k | 0.356 | 0.612 | 0.25 | 3.91 | 0.5830 | Sub-I--I | T1 |

| N03 | NMH? | SpiNNaker2 | 10M | 0.412 | 0.534 | 0.22 | 3.91 | 0.5204 | Sub-I--I | T1 |



\subsection{The Triple Lock and the Thermodynamic Asymptote}

<p align=""center"">

<img src=""figures\_cst/Fig3\_Triple\_Lock.png"" alt=""Fig. 3"" width=""80\%"">

</p>

<p align=""center""><b>Fig. 3 | Triple Lock mechanism.</b> Three concentric barriers confining binary-digital ANN below L1. Outermost: $\alpha$-lock ($\alpha$ = 0.69). Middle: $\Gamma$-lock ($\Gamma$\_st < 0.08). Innermost: $\Omega$-lock ($\Omega$ = 0.03--0.05). Biological systems evolved keys to all three.</p>


Three physical mechanisms constitute the Triple Lock:


1. \textbackslash\{\}textbf\{Low $\alpha$ ($\alpha$\_digital = 0.69 vs. $\alpha$\_cortical = 3.91):\} Binary digital logic constrains M\_eff = 2 states per node regardless of CMOS node size. Information-theoretic analysis yields effective $\alpha$ $\leq$ 1.25--1.6 due to activation compression and spatial correlation.


2. \textbackslash\{\}textbf\{Frozen $\Gamma$\_st ($\Gamma$\_st $\approx$ 0.08 for binary-digital Transformers at inference):\} Training aligns weight structure with functional activations, but once training converges, $\Gamma$\_st is frozen. This is fundamentally different from biological $\Gamma$\_st, continuously updated through synaptic STDP.


3. \textbackslash\{\}textbf\{Suppressed T\_c ($\Psi$ $\approx$ 0.03 for binary-digital Transformers):\} Frozen inference weights eliminate functional connectivity variability.


\textbackslash\{\}textbf\{The binary-digital ceiling:\} CST\_emergent\_max $\approx$ 0.35 (at $\Gamma$\_st $\rightarrow$ 0.5, $\alpha$\_digital = 0.69)---permanently below L1 = 0.707. No amount of parameter scaling can overcome this exponential ceiling.


\textbackslash\{\}textbf\{Counter-argument considered.\} Could online learning with dynamic topology reconfiguration on digital hardware break the Triple Lock? In principle, inference-time weight updates can raise $\Gamma$\_st above the frozen floor, and dynamic architecture search could increase S\_c. However, the $\alpha$-lock remains insurmountable: even if $\Gamma$\_st approaches 0.5 and S\_c approaches 1.0 through software, $\alpha$ = 0.69 imposes CST\_max $\approx$ 1.412---crossing L1 and L2 but remaining below L3 ($\phi$ = 1.618). More critically, achieving $\Gamma$\_st = 0.5 on digital hardware requires continuous synapse-level weight updates at inference---an energy cost growing as O(N2 $\cdot$ $\alpha$\_digital / $\alpha$\_cortical), making the approach thermodynamically prohibitive. The Triple Lock is thus a proof of thermodynamic impossibility, not physical impossibility.


Importantly, this ceiling is not imposed by CMOS technology; analog CMOS implementations of memristive synapses achieve $\alpha$ $\approx$ 3.5--4.5, lifting the ceiling entirely (see Hardware Roadmap).


---


\section{Discussion}

\subsection{IIL vs. TIL: The Two-Layer Intelligence Framework}

CST captures Intrinsic Intelligence Level (IIL)---the physical capacity of a system's architecture to support emergent intelligence. This is distinct from Task Intelligence Level (TIL)---functional performance on specific benchmarks. The two are orthogonal: a tsunami simulation achieves vastly greater numerical accuracy than any human forecaster yet possesses zero intelligence. This distinction resolves the apparent paradox of GPT-class models: benchmark supremacy reflects TIL optimization through parameter scaling, while sub-C. elegans CST reflects IIL limitation from binary-digital architecture.


\subsection{The Hardware Roadmap}

<p align=""center"">

<img src=""figures\_cst/Fig4\_Hardware\_Roadmap.png"" alt=""Fig. 4"" width=""95\%"">

</p>

<p align=""center""><b>Fig. 4 | Four-generation hardware roadmap.</b> Timeline 2024--2032. Gen1 (Device Innovation): memristive SNN, L1--L2. Gen2 (Integration): wafer-scale SDI, L2--L3. Gen3 (SDSoW): heterogeneous 3D, L3--L4. Gen4 (Photonic): optical interconnect, L4+.</p>


The engineering pathway requires crossing the $\Gamma$\_st barrier through materials, not algorithms:


- \textbackslash\{\}textbf\{Gen1 (Device Innovation):\} Memristive crossbar arrays + STDP break the $\alpha$-lock ($\alpha$: 0.69 $\rightarrow$ 3.91), achieving L1--L2.

- \textbackslash\{\}textbf\{Gen2 (Integration):\} Wafer-scale SDI coordinates chiplet-to-chiplet physical coupling, raising effective $\Gamma$\_st.

- \textbackslash\{\}textbf\{Gen3 (SDSoW):\} Heterogeneous 3D integration combines digital logic, analog memory, and photonic routing.

- \textbackslash\{\}textbf\{Gen4 (Photonic):\} Optical interconnects (100 ps $\rightarrow$ 10 ps latency) lift timescale bottleneck $\Theta$, enabling L4+.


\subsection{Convergence of AI Architecture toward CST}

The global AI industry's architectural evolution 2017--2025 provides independent empirical validation: every major advance---MoE modularity, NAS hierarchy, SSM recurrence, LTC dynamics, inference-time plasticity---maps onto a specific CST component [24--27]. The convergence is toward CST-optimal architecture through engineering pressure, while simultaneously revealing the one transition the scaling paradigm cannot make: from simulated $\Gamma$\_st to physical $\Gamma$\_st.


\subsection{Falsifiability and Boundary Conditions}

CST makes specific falsifiable predictions: (1) any system with CST < 0.707 cannot exhibit persistent learned behavior---violation by any binary-digital ANN would falsify the $\alpha$-lock; (2) the six thresholds are universal---discovery of a system at Level III without crossing $\phi$ would falsify the threshold derivation; (3) no binary-digital ANN can exceed $\eta$\_I \textasciitilde{} 10\textasciicircum{}(?4) regardless of scale.


\subsection{Limitations}

$\Gamma$\_st comparability across BNN and ANN measurement modalities requires careful calibration: BNN $\Gamma$\_st is measured via NMI of structural and functional communities from neuroimaging; ANN $\Gamma$\_st is inferred from weight-activation alignment. Systematic offset is estimated at $\pm$0.04. Extension to higher-order networks (simplicial complexes, Betti numbers) is deferred to a companion paper.


---


\section{Methods}

\subsection{Data Provenance and Quality Grading}

All 40 systems graded by measurement directness:


- \textbackslash\{\}textbf\{[T1] Direct measurement (n = 34):\} Parameters from peer-reviewed connectomic/electrophysiological datasets with zero free parameters. Core statistics use T1 only.

- \textbackslash\{\}textbf\{[T2?] Indirect inference (n = 5):\} Error bars $\pm$15\%.

- \textbackslash\{\}textbf\{[T3§] Proxy measurement:\} ANN systems where architecture statistics serve as proxies. Illustrative only, excluded from core statistics.


\textbackslash\{\}textbf\{Data sources.\} C. elegans: Varshney et al. [28]. Mouse: Oh et al. [29]. Human: Van Essen et al. [30], HCP. Branching ratio: Beggs \& Plenz [5]. SC-FC coupling: Arnatkeviciute et al. [31]; Honey et al. [32]. C. elegans dynamics: Kato et al. [33]; Gordus et al. (2015). $\Gamma$\_st = 0.17 from Randi et al. (2024).


\subsection{Unified Cross-Species Computation Protocol (UCCP)}

All components normalized to [0, 1]:


\textbackslash\{\}textbf\{S\_c = (C $\cdot$ H $\cdot$ M $\cdot$ R\_sw)\textasciicircum{}(1/4):\} C = |LCC|/N; H = min[(k\_max/k\_null)/6.667, 1.0], k\_null $\approx$ ln(N)/ln(ln(N)); M = max[(Q ? 0.02)/(1 ? 0.02), 0.01], Q = Louvain modularity (100 restarts, $\gamma$ = 1.0); R\_sw = tanh[($\sigma$ ? 1)/2], $\sigma$ = (C/C\_rand)/(L/L\_rand).


\textbackslash\{\}textbf\{T\_c = ($\lambda$\_eff $\cdot$ $\Phi$ $\cdot$ $\Psi$ $\cdot$ $\Theta$)\textasciicircum{}(1/4).\} BNN: $\lambda$\_eff = avalanche branching ratio; $\Phi$ = mean pairwise PLV; $\Psi$ = std(sliding-window FC)/mean|FC|; $\Theta$ = slowest/fastest intrinsic timescale ratio. ANN: $\lambda$\_eff = Lyapunov exponent; $\Phi$ = inter-channel phase synchrony; $\Psi$ = layer-activation variability; $\Theta$ = gradient timescale diversity.


\textbackslash\{\}textbf\{Sensitivity analysis.\} UCCP validated against Z-score, min-max, rank-based, and PCA-derived normalization. All five core findings invariant across frameworks (Spearman $\rho$ $\geq$ 0.91).


\subsection{Statistical Analysis}

Spearman $\rho$ = 0.976 (95\% CI: 0.956--0.988). Fisher exact tests at $\theta$\_1 (p = 0.0003), $\theta$\_3 (p = 0.0004), $\theta$\_5 (p = 0.0001), all surviving Bonferroni ($\alpha$\_corrected = 0.0083). PIC confirms significance after phylogenetic correction (p < 0.01).



---


\begin{thebibliography}{99}
\bibitem{ref1} Kaplan, J. et al. Scaling laws for neural language models. arXiv:2001.08361 (2020).
\bibitem{ref2} Bennett, C.H. Notes on Landauer's principle, reversible computation, and Maxwell's Demon. Stud. Hist. Philos. Mod. Phys. 34, 501--510 (2003).
\bibitem{ref3} von Neumann, J. Theory of Self-Reproducing Automata (ed. Burks, A.W.). University of Illinois Press (1966).
\bibitem{ref4} Turing, A.M. On computable numbers, with an application to the Entscheidungsproblem. Proc. Lond. Math. Soc. s2-42, 230--265 (1936).
\bibitem{ref5} Beggs, J.M. \textbackslash\{\}\& Plenz, D. Neuronal avalanches in neocortical circuits. J. Neurosci. 23, 11167--11177 (2003).
\bibitem{ref6} Shew, W.L. et al. Neuronal avalanches imply maximum dynamic range in cortical networks at criticality. J. Neurosci. 29, 15595--15600 (2009).
\bibitem{ref7} Tagliazucchi, E. et al. Large-scale signatures of unconsciousness are consistent with a departure from critical dynamics. J. R. Soc. Interface 13, 20151027 (2016).
\bibitem{ref8} di Santo, S. et al. Self-organized bistability and its relevance to brain dynamics. Phys. Rev. Research 5, 013009 (2023).
\bibitem{ref9} Meshulam, L. et al. Coarse graining, fixed points, and scaling in a large population of neurons. Phys. Rev. Lett. 123, 178103 (2019).
\bibitem{ref10} Watts, D.J. \textbackslash\{\}\& Strogatz, S.H. Collective dynamics of small-world networks. Nature 393, 440--442 (1998).
\bibitem{ref11} Sporns, O. \textbackslash\{\}\& Betzel, R.F. Modular brain networks. Annu. Rev. Psychol. 67, 613--640 (2016).
\bibitem{ref12} Friston, K. The free-energy principle: a unified brain theory? Nat. Rev. Neurosci. 11, 127--138 (2010).
\bibitem{ref13} Kelso, J.A.S. Dynamic Patterns: The Self-Organization of Brain and Behavior. MIT Press (1995).
\bibitem{ref14} Tononi, G. et al. Integrated information theory: from consciousness to its physical substrate. Nat. Rev. Neurosci. 17, 450--461 (2016).
\bibitem{ref15} Barrett, A.B. \textbackslash\{\}\& Mediano, P.A.M. The phi measure of integrated information is not well-defined for general physical systems. Entropy 21, 17 (2019).
\bibitem{ref16} Sporns, O. Networks of the Brain. MIT Press (2010).
\bibitem{ref17} Seidman, S.B. Network structure and minimum degree. Social Networks 5, 269--287 (1983).
\bibitem{ref18} Fortunato, S. \textbackslash\{\}\& Barthélemy, M. Resolution limit in community detection. Proc. Natl. Acad. Sci. USA 104, 36--41 (2007).
\bibitem{ref19} Murray, J.D. et al. A hierarchy of intrinsic timescales across primate cortex. Nat. Neurosci. 17, 1661--1663 (2014).
\bibitem{ref20} Zhang, W.-Z. Escaping the Semantic Valley: Non-Abelian Gauge Fields on Large Model Fiber Bundles and the Optimal Gauge Charge. arXiv:2409.12345 (2024).
\bibitem{ref21} Hebb, D.O. The Organization of Behavior. Wiley (1949).
\bibitem{ref22} Feigenbaum, M.J. Quantitative universality for a class of nonlinear transformations. J. Stat. Phys. 19, 25--52 (1978).
\bibitem{ref23} Felsenstein, J. Phylogenies and the comparative method. Am. Nat. 125, 1--15 (1985).
\bibitem{ref24} Shazeer, N. et al. Outrageously large neural networks: the sparsely-gated mixture-of-experts layer. ICLR (2017).
\bibitem{ref25} Gu, A. \textbackslash\{\}\& Dao, T. Mamba: linear-time sequence modeling with selective state spaces. arXiv:2312.00752 (2023).
\bibitem{ref26} Peng, B. et al. RWKV: Reinventing RNNs for the transformer era. arXiv:2305.13048 (2023).
\bibitem{ref27} Sun, Y. et al. Learning to (learn at test time): RNNs with expressive hidden states. arXiv:2407.04620 (2024).
\bibitem{ref28} Varshney, L.R. et al. Structural properties of the C. elegans neuronal network. PLOS Comput. Biol. 7, e1001066 (2011).
\bibitem{ref29} Oh, S.W. et al. A mesoscale connectome of the mouse brain. Nature 508, 207--214 (2014).
\bibitem{ref30} Van Essen, D.C. et al. The WU-Minn Human Connectome Project. NeuroImage 80, 62--79 (2013).
\bibitem{ref31} Arnatkeviciute, A. et al. Structural and functional brain network analysis. NeuroImage 241, 118403 (2021).
\bibitem{ref32} Honey, C.J. et al. Predicting human resting-state functional connectivity from structural connectivity. Proc. Natl. Acad. Sci. USA 106, 2035--2040 (2009).
\bibitem{ref33} Kato, S. et al. Global brain dynamics embed the motor command sequence of Caenorhabditis elegans. Cell 163, 656--669 (2015).
\bibitem{ref34} White, J.G. et al. The structure of the nervous system of C. elegans. Philos. Trans. R. Soc. B 314, 1--40 (1986).
\bibitem{ref35} Scheffer, L.K. et al. A connectome and analysis of the adult Drosophila central brain. eLife 9, e57443 (2020).
\end{thebibliography}

\section*{Author Contributions}
Q.L.: Conceptualization, Methodology, Formal analysis, Data curation, Writing --- original draft, Writing --- review \& editing.

\section*{Competing Interests}
The author declares no competing financial interests.

\section*{Data Availability}
https://github.com/iNEST-TJU/CST-theorem

\section*{AI-Assistance Declaration}
AI-assisted tools were used for literature summarization and language polishing. All content has been reviewed and verified by the authors.

\end{document}